\documentclass[11pt]{article}
\usepackage{graphicx}
\usepackage[margin=1in]{geometry}

\title{A Tiny Paper Built by Make}
\author{Missing Semester --- Metaprogramming exercise}
\date{\today}

\begin{document}
\maketitle

\section{Introduction}
This document exists to demonstrate a real build system. The figure below is
generated from \texttt{data.dat} by \texttt{plot-data.py}; the PDF is produced
from this \texttt{.tex} source. \texttt{make} wires the dependency chain
together so that touching either the data or the plotting script rebuilds only
what is necessary.

\section{Results}
\begin{figure}[h]
  \centering
  \includegraphics[width=0.6\textwidth]{plot.png}
  \caption{Data plotted by the build pipeline.}
\end{figure}

\section{Conclusion}
The build is fully described by the \texttt{Makefile}: \texttt{paper.pdf}
depends on \texttt{plot.png}, which depends on \texttt{data.dat} and the
plotting script. A \texttt{clean} target removes every generated artifact.

\end{document}
